% GENERATED by bin/curate/gen_tutorials_guide.py -- DO NOT EDIT BY HAND.
% Assembled from each case's controlDict description + header comment.

\clearpage
\tutentry{ctrl01\_cstr\_temp\_control}{tutorials/ctrl/ctrl01\_cstr\_temp\_control}
\begin{tutscheme}
\begin{tikzpicture}[>={Stealth[length=2mm]}, every node/.style={font=\scriptsize}]
  \node[draw, rounded corners=2pt, align=center, minimum width=1.9cm, minimum height=0.85cm, fill=black!4] (reactor) at (0.0,-0.0) {\textbf{reactor}\\[-1pt]{\tiny dynamicCSTR}};
  \draw[->] ($(reactor.west)+(-1.05,0.00)$) -- ($(reactor.west)+(0,0.00)$) node[at start, above, font=\tiny, anchor=south west] {feed};
  \draw[->] ($(reactor.east)+(0,0.00)$) -- ($(reactor.east)+(1.05,0.00)$) node[at end, above, font=\tiny, anchor=south east] {product};
\end{tikzpicture}
\end{tutscheme}
\tutwhat{Dynamic CSTR (compA $\to$ compB) with PID jacket-temperature control; T climbs from 320 K into the +/-2\% setpoint band within \textasciitilde{}20 s (no overshoot) and tracks toward 350 K across the 150 s window}
\begin{tutnarrative}
the canonical first choupoCtrl test.

Process: Continuous CSTR with first-order reaction compA $\to$ compB -- an ISOMERISATION: both stand-ins carry MW 30, so the pseudo-reaction conserves mass exactly and the run-end material ledger closes to machine precision. (The element balance stays honestly UNAVAILABLE: compA/compB are numerical stand-ins with no real formula, and the ledger says so by name instead of fabricating atoms.) Feed is pure compA at 320 K, 5$\times$10?? kmol/s; reactor volume 1 L; mild exotherm ($\Delta$H\_rxn = -50 kJ/mol). Without control the reactor settles a few K above the inlet temperature (small reaction exotherm balanced by convective cooling), well below the desired setpoint of 350 K.

Controller: A PID acts on the jacket temperature (T\_jacket) to drive the reactor temperature (T) to a setpoint of 350 K. Tuning is pedagogical, not optimal:

Kp = 4.0 (proportional band $\approx$ 7-8 K) Ki = 0.04 (integral time $\tau$\_I = Kp/Ki $\approx$ 100 s) Kd = 0.0 (no derivative; system is benign enough)

Jacket bounded to [280, 420] K --- realistic limits of a utility loop. Anti-windup: integral does not advance when the unclamped output is outside the bounds.

Expected trajectory (150 s window): Open-loop: T relaxes a few K above the 320 K inlet (small exotherm balanced by convective cooling). Closed-loop: T rises fast (the jacket MV slams to its 420 K clamp, then backs off), entering the +/-2\% setpoint band by \textasciitilde{} 20 s, well-damped with no overshoot. By 150 s T is \textasciitilde{} 348.5 K and still creeping up: the last \textasciitilde{} 1.5 K rides the slow chemical tail (the exotherm's conversion, tau \textasciitilde{} 120 s) toward 350 K beyond the window. The jacket (MV) settles around 348.5 K to keep heating in against the cool 320 K feed. (Widen endTime to watch the final approach to 350 K; the 150 s window is sized to SHOW the control transient, not the slow tail.)

To run the OPEN-LOOP comparison, comment out the entire `controllers` block below --- the case will then run the bare reactor with whatever T\_jacket the `operation` block initialised it to (here 320 K, i.e. no jacket heating).

ONE DECLARED FIGURE (system/postDict, C487): `response`, "how fast does the reactor reach its set point, and how hard does the jacket work?" -- TC1's PV against its set point, and its MV (the jacket temperature, on its 420 K limit at first) on the right-hand axis.

\end{tutnarrative}
\tutreadme{A 1-litre continuous stirred reactor with a first-order A $\to$ B reaction (two stand-in components, both MW 30 --- the pseudo-reaction says so by name rather than inventing atoms), fed at 320 K, with a jacket whose temperature a PID controller sets to hold the reactor at 350 K. 150 s of transient, 1 s steps. The golden: PV ends at 348.47 K, the jacket (MV) at 348.54 K, mass in = out = 0.225 kg over the run and the material balance closes at 1.8e-15.}
\tutreadme{1. The controller is a unit in the dict, with its tuning in the open. type PID; setpoint 350.0; gains \{ Kp 4.0; Ki 0.04; Kd 0.0; \} output \{ min 280; max 420; bias 320; \} --- proportional band $\approx$ 7--8 K, integral time $\tau$\_I = Kp/Ki = 100 s, no derivative. Nothing is tuned for you. 2. Watch the manipulated variable saturate. At t = 2 s the jacket demand jumps to the 420 K clamp (the log's last column), then backs off as the reactor warms: 402, 385, 374 K ? The Log tab prints every 2 s; the clamp is a number in the dict and you can see it bind. 3. Into the band, then the slow part. The reactor crosses 343 K --- the -2 \% edge of the setpoint band --- at t $\approx$ 14 s. At 150 s it sits 1.5 K under setpoint and still climbing: that is the integral action at $\tau$\_I = 100 s working off the offset. Widen endTime and watch it arrive. 4. The energy balance REFUSES here, and says why. [balance] energy: UNAVAILABLE --- the dynamicCSTR energy equation is a Cp/convective model, not the exact derivative of a stored H. The stand-in components carry no formation enthalpy on purpose; a ledger that cannot be exact says so rather than printing a closure it did not earn. ctrl11 is the case where it can, and does (1.4e-7).}
\tutreadme{Set Kp 8.0;: faster into the band, more overshoot, the clamp held longer. Then Ki 0; --- the offset never closes, which is what a proportional-only loop teaches in one run. ctrl02\_disturbance\_rejection steps the feed temperature and asks the same loop to hold.}
\begin{tutgolden}
\goldrow{kpi}{TC1}{MV\_final}{348.54491648}
\goldrow{kpi}{TC1}{PV\_final}{348.468973192}
\goldrow{kpi}{TC1}{SP\_final}{350}
\goldrow{kpi}{reactor}{T\_final}{348.481893564}
\goldrow{kpi}{reactor}{n\_compA\_final}{0.0116161140368}
\goldrow{kpi}{reactor}{n\_compB\_final}{0.00038388596316}
\goldrow{kpi}{reactor}{t\_end}{150}
\goldrow{kpi}{balance}{energy\_balance\_available}{0}
\goldrow{kpi}{balance}{mass\_closure\_rel}{1.77327288655e-15}
\goldrow{kpi}{balance}{mass\_kg\_final}{0.36}
\goldrow{kpi}{balance}{mass\_kg\_in\_cum}{0.225}
\goldrow{kpi}{balance}{mass\_kg\_initial}{0.36}
\goldrow{kpi}{balance}{mass\_kg\_out\_cum}{0.225}
\goldrow{kpi}{balance}{mass\_residual\_kg}{6.38378239159e-16}
\goldmore{4}
\end{tutgolden}

\clearpage
\tutentry{ctrl02\_disturbance\_rejection}{tutorials/ctrl/ctrl02\_disturbance\_rejection}
\begin{tutscheme}
\begin{tikzpicture}[>={Stealth[length=2mm]}, every node/.style={font=\scriptsize}]
  \node[draw, rounded corners=2pt, align=center, minimum width=1.9cm, minimum height=0.85cm, fill=black!4] (reactor) at (0.0,-0.0) {\textbf{reactor}\\[-1pt]{\tiny dynamicCSTR}};
  \draw[->] ($(reactor.west)+(-1.05,0.00)$) -- ($(reactor.west)+(0,0.00)$) node[at start, above, font=\tiny, anchor=south west] {feed};
  \draw[->] ($(reactor.east)+(0,0.00)$) -- ($(reactor.east)+(1.05,0.00)$) node[at end, above, font=\tiny, anchor=south east] {product};
\end{tikzpicture}
\end{tutscheme}
\tutwhat{PID disturbance rejection: T\_in steps down then up, PID maintains T\_setpoint}
\begin{tutnarrative}
the textbook control demo.

Same DynamicCSTR as ctrl01 (compA $\to$ compB, mild exotherm, jacket heat transfer). Same PID temperature controller. Two added ingredients:

1. The simulation runs 2000 s (long enough to reach steady state after each disturbance). 2. A second "controller" of type Schedule injects step disturbances in the inlet temperature T\_in:

0 | 320 (nominal) 700 | 305 (-15 K cold disturbance) 1300 | 335 (+15 K warm disturbance, back-and-up)

The PID controller has no foreknowledge of these disturbances --- it only sees T\_reactor drift away from setpoint and acts on the jacket to compensate. Expected behaviour:

0 < t < 700: normal startup, T $\to$ 350 K ($\approx$ ctrl01). t = 700: T\_in drops 15 K; the integral action absorbs it so tightly that T\_reactor holds at 350 K with no visible sag (the jacket simply keeps heating). t = 1300: T\_in rises 15 K; again T\_reactor stays pinned at 350 K with no visible deviation.

The trajectory.csv columns now include both controllers' SP / PV / MV. The Schedule controller's SP and MV mirror the same step train (it has no PV); the PID's SP stays at 350 K and its PV (T\_reactor) tracks closely.

TWO DECLARED FIGURES (system/postDict, C487): `response`, "does the controller hold the reactor at its set point?" -- the controlled variable (TC1's PV) against its set point, the manipulated variable (TC1's MV, the jacket temperature) and the disturbance (FeedDist's MV, the feed temperature), all against time; and `rejection`, the same from 600 s on a tight temperature axis, where the dip and peak after the two steps can be read. Both are drawn from the `trajectory` table each controller publishes.

\end{tutnarrative}
\begin{tutgolden}
\goldrow{kpi}{FeedDist}{MV\_final}{335}
\goldrow{kpi}{TC1}{MV\_final}{347.976535585}
\goldrow{kpi}{TC1}{PV\_final}{350.00207645}
\goldrow{kpi}{TC1}{SP\_final}{350}
\goldrow{kpi}{reactor}{T\_final}{350.002059115}
\goldrow{kpi}{reactor}{n\_compA\_final}{0.0110806465498}
\goldrow{kpi}{reactor}{n\_compB\_final}{0.000919353450202}
\goldrow{kpi}{reactor}{t\_end}{2000}
\goldrow{kpi}{balance}{energy\_balance\_available}{0}
\goldrow{kpi}{balance}{mass\_closure\_rel}{2.79221090693e-14}
\goldrow{kpi}{balance}{mass\_kg\_final}{0.36}
\goldrow{kpi}{balance}{mass\_kg\_in\_cum}{3}
\goldrow{kpi}{balance}{mass\_kg\_initial}{0.36}
\goldrow{kpi}{balance}{mass\_kg\_out\_cum}{3}
\goldmore{5}
\end{tutgolden}

\clearpage
\tutentry{ctrl03\_adaptive\_disturbance}{tutorials/ctrl/ctrl03\_adaptive\_disturbance}
\begin{tutscheme}
\begin{tikzpicture}[>={Stealth[length=2mm]}, every node/.style={font=\scriptsize}]
  \node[draw, rounded corners=2pt, align=center, minimum width=1.9cm, minimum height=0.85cm, fill=black!4] (reactor) at (0.0,-0.0) {\textbf{reactor}\\[-1pt]{\tiny dynamicCSTR}};
  \draw[->] ($(reactor.west)+(-1.05,0.00)$) -- ($(reactor.west)+(0,0.00)$) node[at start, above, font=\tiny, anchor=south west] {feed};
  \draw[->] ($(reactor.east)+(0,0.00)$) -- ($(reactor.east)+(1.05,0.00)$) node[at end, above, font=\tiny, anchor=south east] {product};
\end{tikzpicture}
\end{tutscheme}
\tutwhat{PID disturbance rejection with ADAPTIVE plant integration --- controller samples on the fixed 1 s grid, Rosenbrock23 sub-steps between samples (MV held)}
\begin{tutnarrative}
the ctrl02\_disturbance\_rejection demo re-integrated with an ADAPTIVE plant integrator: the PID still samples on the fixed 1 s control grid (MV held between samples) while Rosenbrock23 sub-steps the plant adaptively in between. The flowsheet below is ctrl02's textbook disturbance-rejection setup, unchanged:

Same DynamicCSTR as ctrl01 (compA $\to$ compB, mild exotherm, jacket heat transfer). Same PID temperature controller. Two added ingredients:

1. The simulation runs 2000 s (long enough to reach steady state after each disturbance). 2. A second "controller" of type Schedule injects step disturbances in the inlet temperature T\_in:

0 | 320 (nominal) 700 | 305 (-15 K cold disturbance) 1300 | 335 (+15 K warm disturbance, back-and-up)

The PID controller has no foreknowledge of these disturbances --- it only sees T\_reactor drift away from setpoint and acts on the jacket to compensate. Expected behaviour:

0 < t < 700: normal startup, T $\to$ 350 K ($\approx$ ctrl01). t = 700: T\_in drops 15 K; the integral action absorbs it so tightly that T\_reactor holds at 350 K with no visible sag (the jacket simply keeps heating). t = 1300: T\_in rises 15 K; again T\_reactor stays pinned at 350 K with no visible deviation.

The trajectory.csv columns now include both controllers' SP / PV / MV. The Schedule controller's SP and MV mirror the same step train (it has no PV); the PID's SP stays at 350 K and its PV (T\_reactor) tracks closely.

TWO DECLARED FIGURES (system/postDict, C487), ctrl02's: `response`, "does the controller hold the reactor at its set point?", and `rejection`, the same from 600 s on a tight temperature axis (349.5 to 350.8 K), where the 0.35 K dip and 0.66 K peak after the two feed steps can be read.

\end{tutnarrative}
\begin{tutgolden}
\goldrow{kpi}{FeedDist}{MV\_final}{335}
\goldrow{kpi}{TC1}{MV\_final}{347.976535626}
\goldrow{kpi}{TC1}{PV\_final}{350.002076468}
\goldrow{kpi}{TC1}{SP\_final}{350}
\goldrow{kpi}{reactor}{T\_final}{350.002059134}
\goldrow{kpi}{reactor}{n\_compA\_final}{0.0110806465416}
\goldrow{kpi}{reactor}{n\_compB\_final}{0.000919353458409}
\goldrow{kpi}{reactor}{t\_end}{2000}
\goldrow{kpi}{balance}{energy\_balance\_available}{0}
\goldrow{kpi}{balance}{mass\_closure\_rel}{2.66083451568e-14}
\goldrow{kpi}{balance}{mass\_kg\_final}{0.36}
\goldrow{kpi}{balance}{mass\_kg\_in\_cum}{3}
\goldrow{kpi}{balance}{mass\_kg\_initial}{0.36}
\goldrow{kpi}{balance}{mass\_kg\_out\_cum}{3}
\goldmore{5}
\end{tutgolden}

\clearpage
\tutentry{ctrl05\_disturbance\_transient}{tutorials/ctrl/ctrl05\_disturbance\_transient}
\begin{tutscheme}
\begin{tikzpicture}[>={Stealth[length=2mm]}, every node/.style={font=\scriptsize}]
  \node[draw, rounded corners=2pt, align=center, minimum width=1.9cm, minimum height=0.85cm, fill=black!4] (reactor) at (0.0,-0.0) {\textbf{reactor}\\[-1pt]{\tiny dynamicCSTR}};
  \draw[->] ($(reactor.west)+(-1.05,0.00)$) -- ($(reactor.west)+(0,0.00)$) node[at start, above, font=\tiny, anchor=south west] {feed};
  \draw[->] ($(reactor.east)+(0,0.00)$) -- ($(reactor.east)+(1.05,0.00)$) node[at end, above, font=\tiny, anchor=south east] {product};
\end{tikzpicture}
\end{tutscheme}
\tutwhat{PID disturbance rejection with OpenFOAM-style real-time instant dirs: T\_in steps down (700 s) then up (1300 s); scrub 0/ 40/ 80/ ... s to watch the PID track the 350 K setpoint and the jacket MV move}
\begin{tutnarrative}
the textbook control demo, captured in OpenFOAM-style real-time instant directories (see controlDict's solutionControl block). Identical physics to ctrl02; the difference is that each written time directory (the layout of 0/: the holdup under `<t>/internalStates/`, one file per stream) can be scrubbed across the two disturbance steps to watch the PID reject them.

Same DynamicCSTR as ctrl01 (compA $\to$ compB, mild exotherm, jacket heat transfer). Same PID temperature controller. Two added ingredients:

1. The simulation runs 2000 s (long enough to reach steady state after each disturbance). 2. A second "controller" of type Schedule injects step disturbances in the inlet temperature T\_in:

0 | 320 (nominal) 700 | 305 (-15 K cold disturbance) 1300 | 335 (+15 K warm disturbance, back-and-up)

The PID controller has no foreknowledge of these disturbances --- it only sees T\_reactor drift away from setpoint and acts on the jacket to compensate. Expected behaviour:

0 < t < 700: normal startup, T $\to$ 350 K ($\approx$ ctrl01). t = 700: T\_in drops 15 K; the integral action absorbs it so tightly that T\_reactor holds at 350 K with no visible sag (the jacket simply keeps heating). t = 1300: T\_in rises 15 K; again T\_reactor stays pinned at 350 K with no visible deviation.

The trajectory.csv columns now include both controllers' SP / PV / MV. The Schedule controller's SP and MV mirror the same step train (it has no PV); the PID's SP stays at 350 K and its PV (T\_reactor) tracks closely.

TWO DECLARED FIGURES (system/postDict, C487), ctrl02's: `response`, "does the controller hold the reactor at its set point?", and `rejection`, the same from 600 s on a tight temperature axis (349.5 to 350.8 K), where the 0.35 K dip and 0.66 K peak after the two feed steps can be read.

\end{tutnarrative}
\begin{tutgolden}
\goldrow{kpi}{FeedDist}{MV\_final}{335}
\goldrow{kpi}{TC1}{MV\_final}{347.976535585}
\goldrow{kpi}{TC1}{PV\_final}{350.00207645}
\goldrow{kpi}{TC1}{SP\_final}{350}
\goldrow{kpi}{reactor}{T\_final}{350.002059115}
\goldrow{kpi}{reactor}{n\_compA\_final}{0.0110806465498}
\goldrow{kpi}{reactor}{n\_compB\_final}{0.000919353450202}
\goldrow{kpi}{reactor}{t\_end}{2000}
\goldrow{kpi}{balance}{energy\_balance\_available}{0}
\goldrow{kpi}{balance}{mass\_closure\_rel}{2.79221090693e-14}
\goldrow{kpi}{balance}{mass\_kg\_final}{0.36}
\goldrow{kpi}{balance}{mass\_kg\_in\_cum}{3}
\goldrow{kpi}{balance}{mass\_kg\_initial}{0.36}
\goldrow{kpi}{balance}{mass\_kg\_out\_cum}{3}
\goldmore{5}
\end{tutgolden}

\clearpage
\tutentry{ctrl06\_sine\_disturbance}{tutorials/ctrl/ctrl06\_sine\_disturbance}
\begin{tutscheme}
\begin{tikzpicture}[>={Stealth[length=2mm]}, every node/.style={font=\scriptsize}]
  \node[draw, rounded corners=2pt, align=center, minimum width=1.9cm, minimum height=0.85cm, fill=black!4] (reactor) at (0.0,-0.0) {\textbf{reactor}\\[-1pt]{\tiny dynamicCSTR}};
  \draw[->] ($(reactor.west)+(-1.05,0.00)$) -- ($(reactor.west)+(0,0.00)$) node[at start, above, font=\tiny, anchor=south west] {feed};
  \draw[->] ($(reactor.east)+(0,0.00)$) -- ($(reactor.east)+(1.05,0.00)$) node[at end, above, font=\tiny, anchor=south east] {product};
\end{tikzpicture}
\end{tutscheme}
\tutwhat{Forced SINUSOIDAL inlet-T disturbance: T\_in = 320 + 15 sin(2 pi t / 600). The PID chases a moving target $\to$ the reactor temperature settles into a forced sinusoidal response (attenuated + phase-lagged). Sweep the period to build a one-point Bode by trial-and-error (the frequency-response seed).}
\begin{tutnarrative}
the FREQUENCY-RESPONSE seed.

Same DynamicCSTR as ctrl01/ctrl02 (compA $\to$ compB, mild exotherm, jacket heat transfer) and the same PID temperature controller. The NEW ingredient is the disturbance: instead of the step train of ctrl02's `Schedule` controller, a `Signal` controller injects a SINUSOID into the inlet temperature T\_in:

T\_in(t) = 320 + 15 $\ast$ sin(2$\ast$pi$\ast$t / 600) [K]

mean = 320 K amplitude = 15 K period = 600 s (f = 1/600 = 1.667e-3 Hz, omega = 0.01047 rad/s)

The PID controller has no foreknowledge of the forcing --- it only sees T\_reactor drift sinusoidally away from the 350 K setpoint and works the jacket to compensate. After the first \textasciitilde{}period the reactor temperature settles into a FORCED SINUSOIDAL RESPONSE: same period as the input, but ATTENUATED (the closed loop rejects most of the swing) and PHASE-LAGGED.

THE LESSON (Seborg/Edgar/Mellichamp/Doyle, Ch. 13-14 frequency response): read the output/input amplitude ratio [dB] and the phase lag [deg] off the steady-cycling trajectory, and you have ONE POINT of a Bode plot. Sweep the `period` (edit one number below), re-run, and you trace the closed loop's frequency response by trial-and-error --- the differentiator.

Backward-compatible by construction: ctrl02's `type Schedule; schedule(...)` staircase still works unchanged; this case swaps it for the more general `type Signal; signal \{ type sinusoidal; ... \}`.

TWO DECLARED FIGURES (system/postDict, C487): `response`, "does the reactor stay at its set point while the feed temperature swings?", and `attenuation`, "how much of the 15 K swing of the feed reaches the reactor?" -- from 600 s on a tight axis (349.5 to 350.5 K) beside the feed temperature.

\end{tutnarrative}
\begin{tutgolden}
\goldrow{kpi}{FeedDist}{MV\_final}{319.842923238}
\goldrow{kpi}{FeedDist}{PV\_final}{0}
\goldrow{kpi}{FeedDist}{SP\_final}{319.842923238}
\goldrow{kpi}{TC1}{MV\_final}{350.026347162}
\goldrow{kpi}{TC1}{PV\_final}{350.170890591}
\goldrow{kpi}{TC1}{SP\_final}{350}
\goldrow{kpi}{reactor}{T\_final}{350.173436842}
\goldrow{kpi}{reactor}{n\_compA\_final}{0.0110844237702}
\goldrow{kpi}{reactor}{n\_compB\_final}{0.000915576229809}
\goldrow{kpi}{reactor}{t\_end}{3000}
\goldrow{kpi}{balance}{energy\_balance\_available}{0}
\goldrow{kpi}{balance}{mass\_closure\_rel}{5.72875080707e-14}
\goldrow{kpi}{balance}{mass\_kg\_final}{0.36}
\goldrow{kpi}{balance}{mass\_kg\_in\_cum}{4.5}
\goldmore{7}
\end{tutgolden}

\clearpage
\tutentry{ctrl07\_lhhw\_inhibition}{tutorials/ctrl/ctrl07\_lhhw\_inhibition}
\begin{tutscheme}
\begin{tikzpicture}[>={Stealth[length=2mm]}, every node/.style={font=\scriptsize}]
  \node[draw, rounded corners=2pt, align=center, minimum width=1.9cm, minimum height=0.85cm, fill=black!4] (reactor) at (0.0,-0.0) {\textbf{reactor}\\[-1pt]{\tiny dynamicCSTR}};
  \draw[->] ($(reactor.west)+(-1.05,0.00)$) -- ($(reactor.west)+(0,0.00)$) node[at start, above, font=\tiny, anchor=south west] {feed};
  \draw[->] ($(reactor.east)+(0,0.00)$) -- ($(reactor.east)+(1.05,0.00)$) node[at end, above, font=\tiny, anchor=south east] {product};
\end{tikzpicture}
\end{tutscheme}
\tutwhat{Product inhibition meets a PID: an LHHW reaction whose exotherm fades as its own product covers the catalyst}
\begin{tutnarrative}
ctrl01's temperature loop, with a reaction that poisons itself.

Same vessel, same PID, same setpoint. The one change is the rate law: the product compB now ADSORBS on the catalyst, so the denominator grows as the reaction proceeds,

r = k(T) . c\_A / (1 + K\_B c\_B)

and the reaction throttles itself. A power law cannot express that at any order; only the adsorption term can. This is the same `RateLaw` that the steady `cstr` and the `batchReactor` evaluate -- one definition, three reactors.

What the controller sees: an exotherm that FADES as conversion proceeds, not a constant one. The heat the loop is rejecting shrinks under it, and the jacket temperature must follow. A tuning that holds against a first-order exotherm is not the same tuning that holds against a decaying one.

The adsorption constant K0 = 2.0 m\textasciicircum{}3/kmol is DIDACTIC -- chosen so the inhibition is visible over the control horizon, not measured on any catalyst. For a real, cited LHHW see `cstr07\_lhhw\_methylAcetate` and `batch07\_lhhw\_methylAcetate` (Popken 2000, Eq 16, on UNIQUAC activities).

ONE DECLARED FIGURE (system/postDict, C487): `response`, "can the PID still bring the reactor to its set point when the reaction poisons itself?" -- PV against the set point and the jacket temperature, to be read beside ctrl01's.

\end{tutnarrative}
\begin{tutgolden}
\goldrow{kpi}{TC1}{MV\_final}{349.740705746}
\goldrow{kpi}{TC1}{PV\_final}{348.328939265}
\goldrow{kpi}{TC1}{SP\_final}{350}
\goldrow{kpi}{reactor}{T\_final}{348.342145073}
\goldrow{kpi}{reactor}{n\_compA\_final}{0.0117152677233}
\goldrow{kpi}{reactor}{n\_compB\_final}{0.000284732276726}
\goldrow{kpi}{reactor}{t\_end}{150}
\goldrow{kpi}{balance}{energy\_balance\_available}{0}
\goldrow{kpi}{balance}{mass\_closure\_rel}{7.70988211545e-16}
\goldrow{kpi}{balance}{mass\_kg\_final}{0.36}
\goldrow{kpi}{balance}{mass\_kg\_in\_cum}{0.225}
\goldrow{kpi}{balance}{mass\_kg\_initial}{0.36}
\goldrow{kpi}{balance}{mass\_kg\_out\_cum}{0.225}
\goldrow{kpi}{balance}{mass\_residual\_kg}{2.77555756156e-16}
\goldmore{4}
\end{tutgolden}

\clearpage
\tutentry{ctrl08\_prbs\_ident}{tutorials/ctrl/ctrl08\_prbs\_ident}
\begin{tutscheme}
\begin{tikzpicture}[>={Stealth[length=2mm]}, every node/.style={font=\scriptsize}]
  \node[draw, rounded corners=2pt, align=center, minimum width=1.9cm, minimum height=0.85cm, fill=black!4] (reactor) at (0.0,-0.0) {\textbf{reactor}\\[-1pt]{\tiny dynamicCSTR}};
  \draw[->] ($(reactor.west)+(-1.05,0.00)$) -- ($(reactor.west)+(0,0.00)$) node[at start, above, font=\tiny, anchor=south west] {feed};
  \draw[->] ($(reactor.east)+(0,0.00)$) -- ($(reactor.east)+(1.05,0.00)$) node[at end, above, font=\tiny, anchor=south east] {product};
\end{tikzpicture}
\end{tutscheme}
\tutwhat{PRBS identification experiment: T\_in = 320 +/- 10 K by a maximal-length LFSR (n=4, period 15 bits x 120 s) -- deterministic, seed-declared, period-VERIFIED; the broadband excitation the sine of ctrl06 cannot give.}
\begin{tutnarrative}
PRBS identification experiment (forum \#100.4/\#101).

The classic system-identification input: a pseudo-random binary sequence on the inlet temperature, T\_in = 320 +/- 10 K, driven by a MAXIMAL-LENGTH Fibonacci LFSR --- n = 4 registers, canonical taps (4 3), DECLARED seed 1, period 2\textasciicircum{}4 - 1 = 15 bits of 120 s each (1800 s per full period).

Why a PRBS and not ctrl06's sine: one sinusoid excites ONE frequency; the PRBS is broadband (its autocorrelation approximates an impulse), so a single record carries the information for a full low-order model fit. And unlike real "random" noise it is DETERMINISTIC: seed + taps declared in this dict reproduce the experiment bit for bit, run after run.

THE HAND-DERIVED SEQUENCE (the exact-sequence contract of \#101): with the frozen convention (output = s\_n; feedback = XOR of taps; shift s\_1 <- feedback; seed LSB $\to$ s\_1), n=4 / taps (4 3) / seed 1 gives

0 0 0 1 0 0 1 1 0 1 0 1 1 1 1 (bits k = 0..14, then repeats)

-- walk it yourself: state 0001 $\to$ 0010 $\to$ 0100 $\to$ 1001 $\to$ 0011 $\to$ 0110 $\to$ 1101 $\to$ 1010 $\to$ 0101 $\to$ 1011 $\to$ 0111 $\to$ 1111 $\to$ 1110 $\to$ 1100 $\to$ 1000 $\to$ 0001 (back at the seed after EXACTLY 15 steps: maximal). The run's announce line prints these same first bits; the initialiser REFUSES any tap set whose cycle is not exactly 2\textasciicircum{}n - 1.

The PID (same tuning as ctrl06) holds the reactor at 350 K against the binary disturbance; the trajectory's T\_in column IS the PRBS, and the reactor response is the identification record.

TWO DECLARED FIGURES (system/postDict, C487): `response`, "does the reactor stay at its set point while the feed temperature switches at random?", and `record`, the identification record -- from 600 s the reactor temperature on a tight axis (349.5 to 350.5 K) beside the binary feed temperature.

\end{tutnarrative}
\begin{tutgolden}
\goldrow{kpi}{FeedPRBS}{MV\_final}{330}
\goldrow{kpi}{FeedPRBS}{PV\_final}{0}
\goldrow{kpi}{FeedPRBS}{SP\_final}{330}
\goldrow{kpi}{TC1}{MV\_final}{348.579621229}
\goldrow{kpi}{TC1}{PV\_final}{350.006520504}
\goldrow{kpi}{TC1}{SP\_final}{350}
\goldrow{kpi}{reactor}{T\_final}{350.006465191}
\goldrow{kpi}{reactor}{n\_compA\_final}{0.0110804887195}
\goldrow{kpi}{reactor}{n\_compB\_final}{0.000919511280459}
\goldrow{kpi}{reactor}{t\_end}{3600}
\goldrow{kpi}{balance}{energy\_balance\_available}{0}
\goldrow{kpi}{balance}{mass\_closure\_rel}{6.73638100034e-14}
\goldrow{kpi}{balance}{mass\_kg\_final}{0.36}
\goldrow{kpi}{balance}{mass\_kg\_in\_cum}{5.4}
\goldmore{7}
\end{tutgolden}

\clearpage
\tutentry{ctrl09\_stream\_disturbance}{tutorials/ctrl/ctrl09\_stream\_disturbance}
\begin{tutscheme}
\begin{tikzpicture}[>={Stealth[length=2mm]}, every node/.style={font=\scriptsize}]
  \node[draw, rounded corners=2pt, align=center, minimum width=1.9cm, minimum height=0.85cm, fill=black!4] (reactor) at (0.0,-0.0) {\textbf{reactor}\\[-1pt]{\tiny dynamicCSTR}};
  \draw[->] ($(reactor.west)+(-1.05,0.00)$) -- ($(reactor.west)+(0,0.00)$) node[at start, above, font=\tiny, anchor=south west] {feed};
  \draw[->] ($(reactor.east)+(0,0.00)$) -- ($(reactor.east)+(1.05,0.00)$) node[at end, above, font=\tiny, anchor=south east] {product};
\end{tikzpicture}
\end{tutscheme}
\tutwhat{Inlet-field disturbance: componentMolarFlow.compA steps 3.5e-5 $\to$ 5.25e-5 kmol/s; F and z\_in change by the DERIVED values (F = 6.75e-5, z\_A = 7/9) -- the disturbance never hides from the other components}
\begin{tutnarrative}
the inlet-field actuator, end to end.

The Signal controller targets the reactor's INLET FACE by field name (`inletField componentMolarFlow.compA`) --- the only legal disturbance surface (an outlet actuator is not even expressible). The AUTHORITATIVE inlet state is the per-component molar-flow vector; F\_in and z\_in are DERIVED from it on every write, so a composition disturbance can never hide inside a silent renormalisation (forum \#101):

t < 600 s: nA = 3.5e-5, nB = 1.5e-5 kmol/s $\to$ F\_in = 5.0e-5, z\_A = 0.70, z\_B = 0.30 t >= 600 s: nA stepped to 5.25e-5; nB UNTOUCHED at 1.5e-5 $\to$ F\_in = 6.75e-5, z\_A = 5.25/6.75 = 7/9 = 0.77778, z\_B = 1.5/6.75 = 2/9 = 0.22222 (hand-derived)

The golden pins F\_in and z\_in\_$\ast$ finals to these fractions, and the Signal value is the ABSOLUTE molar flow (announced at start-up). The reactor's own response (more A in, higher conversion load, T transient against the PID) is the pedagogy; the CONTRACT is the derived pair.

The `moleFraction.<c>` sibling field instead holds F\_in constant and renormalises the OTHER components proportionally --- two explicit semantics, each announced, never a guess.

TWO DECLARED FIGURES (system/postDict, C487): `response`, "does the reactor stay at its set point when the feed of compA steps up?" (the compA flow on the right-hand axis), and `step`, the same from 550 s with PV, set point and jacket temperature on one tight axis (349.6 to 350.5 K).

\end{tutnarrative}
\begin{tutgolden}
\goldrow{kpi}{FeedComp}{MV\_final}{5.25e-05}
\goldrow{kpi}{FeedComp}{PV\_final}{0}
\goldrow{kpi}{FeedComp}{SP\_final}{5.25e-05}
\goldrow{kpi}{TC1}{MV\_final}{350.204904403}
\goldrow{kpi}{TC1}{PV\_final}{350.000124372}
\goldrow{kpi}{TC1}{SP\_final}{350}
\goldrow{kpi}{reactor}{F\_in\_final}{6.75e-05}
\goldrow{kpi}{reactor}{T\_final}{350.000123673}
\goldrow{kpi}{reactor}{n\_compA\_final}{0.00879292389459}
\goldrow{kpi}{reactor}{n\_compB\_final}{0.00320707610541}
\goldrow{kpi}{reactor}{t\_end}{1800}
\goldrow{kpi}{reactor}{z\_in\_compA\_final}{0.777777777778}
\goldrow{kpi}{reactor}{z\_in\_compB\_final}{0.222222222222}
\goldrow{kpi}{balance}{energy\_balance\_available}{0}
\goldmore{7}
\end{tutgolden}

\clearpage
\tutentry{ctrl10\_brine\_concentration}{tutorials/ctrl/ctrl10\_brine\_concentration}
\begin{tutscheme}
\begin{tikzpicture}[>={Stealth[length=2mm]}, every node/.style={font=\scriptsize}]
  \node[draw, rounded corners=2pt, align=center, minimum width=1.9cm, minimum height=0.85cm, fill=black!4] (tank) at (0.0,-0.0) {\textbf{tank}\\[-1pt]{\tiny dynamicCSTR}};
  \draw[->] ($(tank.west)+(-1.05,0.00)$) -- ($(tank.west)+(0,0.00)$) node[at start, above, font=\tiny, anchor=south west] {feed};
  \draw[->] ($(tank.east)+(0,0.00)$) -- ($(tank.east)+(1.05,0.00)$) node[at end, above, font=\tiny, anchor=south east] {product};
\end{tikzpicture}
\end{tutscheme}
\tutwhat{Dynamic brine tank on a Pitzer package: the electrolyte thermodynamics through a transient}
\begin{tutnarrative}
ONE stirred vessel, NO reaction. The dynamic CSTR tolerates an empty reaction list, and that is exactly the model wanted here: the holdup is a mixing lag, and the thermodynamics -- Pitzer activity coefficients on the molality basis -- is what travels through it.

The jacket loop is the same PID shape as ctrl01: it acts on T\_jacket to hold the vessel temperature. A brine tank held on temperature is what sits upstream of every evaporator and crystalliser in the corpus.

ONE DECLARED FIGURE (system/postDict, C487): `response`, "does the jacket bring the brine tank to its set point, and does it run out of jacket?" -- PV against the set point, the jacket temperature on its 360 K limit on the right-hand axis.

\end{tutnarrative}
\begin{tutgolden}
\goldrow{kpi}{TC1}{MV\_final}{360}
\goldrow{kpi}{TC1}{PV\_final}{310.319973652}
\goldrow{kpi}{TC1}{SP\_final}{313.15}
\goldrow{kpi}{balance}{element\_Cl\_closure\_rel}{7.17019034075e-16}
\goldrow{kpi}{balance}{element\_H\_closure\_rel}{3.40011338976e-16}
\goldrow{kpi}{balance}{element\_Na\_closure\_rel}{7.17019034075e-16}
\goldrow{kpi}{balance}{element\_O\_closure\_rel}{3.40011338976e-16}
\goldrow{kpi}{balance}{element\_worst\_closure\_rel}{7.17019034075e-16}
\goldrow{kpi}{balance}{energy\_balance\_available}{1}
\goldrow{kpi}{balance}{energy\_closure\_rel}{7.85289304221e-11}
\goldrow{kpi}{balance}{energy\_functional\_H}{1}
\goldrow{kpi}{balance}{energy\_kJ\_final}{-853632.741074}
\goldrow{kpi}{balance}{energy\_kJ\_heat\_cum}{2658.78680029}
\goldrow{kpi}{balance}{energy\_kJ\_in\_cum}{-87796.8114225}
\goldmore{17}
\end{tutgolden}

\clearpage
\tutentry{ctrl11\_esterification\_jacket}{tutorials/ctrl/ctrl11\_esterification\_jacket}
\begin{tutscheme}
\begin{tikzpicture}[>={Stealth[length=2mm]}, every node/.style={font=\scriptsize}]
  \node[draw, rounded corners=2pt, align=center, minimum width=1.9cm, minimum height=0.85cm, fill=black!4] (reactor) at (0.0,-0.0) {\textbf{reactor}\\[-1pt]{\tiny dynamicCSTR}};
  \draw[->] ($(reactor.west)+(-1.05,0.00)$) -- ($(reactor.west)+(0,0.00)$) node[at start, above, font=\tiny, anchor=south west] {feed};
  \draw[->] ($(reactor.east)+(0,0.00)$) -- ($(reactor.east)+(1.05,0.00)$) node[at end, above, font=\tiny, anchor=south east] {product};
\end{tikzpicture}
\end{tutscheme}
\tutwhat{Jacketed esterification CSTR under PID temperature control: the first ctrl case whose stored enthalpy exists, so the first-law ledger claims closure instead of refusing}
\begin{tutnarrative}
One jacketed vessel and one PID loop -- the same shape as ctrl01, so the two can be read side by side. What differs is the chemistry: real species with a formation datum instead of numerical stand-ins, which is the whole reason this case can carry an energy claim and ctrl01 cannot.

The reaction declares NO dH\_rxn. On the elements datum the heat of reaction is Sum(nu\_i h\_i(T)) -- a consequence of the same enthalpy surface the vessel stores, not a second number to keep in step with it.

TWO DECLARED FIGURES (system/postDict, C487): `response`, "how fast does the jacket loop bring the esterification to its set point?", and `holdup`, "how far has the esterification gone?" -- the reactants' and products' amounts in the vessel.

\end{tutnarrative}
\begin{tutgolden}
\goldrow{kpi}{TC1}{MV\_final}{366.642237189}
\goldrow{kpi}{TC1}{PV\_final}{364.93860506}
\goldrow{kpi}{TC1}{SP\_final}{365}
\goldrow{kpi}{balance}{element\_C\_closure\_rel}{1.32609972386e-14}
\goldrow{kpi}{balance}{element\_H\_closure\_rel}{1.31993181817e-14}
\goldrow{kpi}{balance}{element\_O\_closure\_rel}{1.32609972386e-14}
\goldrow{kpi}{balance}{element\_worst\_closure\_rel}{1.32609972386e-14}
\goldrow{kpi}{balance}{energy\_balance\_available}{1}
\goldrow{kpi}{balance}{energy\_closure\_rel}{1.59809574171e-07}
\goldrow{kpi}{balance}{energy\_functional\_H}{1}
\goldrow{kpi}{balance}{energy\_kJ\_final}{-28408.9917781}
\goldrow{kpi}{balance}{energy\_kJ\_heat\_cum}{587.253412702}
\goldrow{kpi}{balance}{energy\_kJ\_in\_cum}{-67689.3808709}
\goldrow{kpi}{balance}{energy\_kJ\_initial}{-28597.6577905}
\goldmore{20}
\end{tutgolden}

\clearpage
\tutentry{ctrl12\_williams\_otto}{tutorials/ctrl/ctrl12\_williams\_otto}
\begin{tutscheme}
\begin{tikzpicture}[>={Stealth[length=2mm]}, every node/.style={font=\scriptsize}]
  \node[draw, rounded corners=2pt, align=center, minimum width=1.9cm, minimum height=0.85cm, fill=black!4] (plant) at (0.0,-0.0) {\textbf{plant}\\[-1pt]{\tiny williamsOttoPlant}};
\end{tikzpicture}
\end{tutscheme}
\tutwhat{Williams-Otto plant (arXiv:2004.07614 eqs. 3.6-3.11 verbatim) settling onto the published steady state x$\ast$ = (3.27, 7.47, 1.12, 9.81, 1.69, 0.22) klb}
\begin{tutnarrative}
One unit, because the benchmark IS one equation set: the Williams-Otto reactor + decanter + column + splitter fold into six ODEs (the recycle term (1-eta)mu, the removed G, the 0.1 F\_tE column retention are all inside them -- see the unit's own header). The five controls sit in `operation`, declared SI: with the MW 453.592 convention (1 kmol == 1 klb) every numeral below IS the paper's value in klb/h.

u$\ast$ of eq. (3.20): F\_fA = 10, F\_fB = 20, T = 580 degR, mu = 129.5, eta = 0.2. T rides in 0/internalStates/plant (322.2222 K = 580 degR exactly).

TWO DECLARED FIGURES (system/postDict, C487): `rates`, "how long does the open-loop plant take to settle at the operating point, and what does it make?" (product and waste rates, klb/h, against hours), and `harvest`, the cumulative yield, waste and objective J (klb).

\end{tutnarrative}
\begin{tutgolden}
\goldrow{kpi}{balance}{energy\_balance\_available}{0}
\goldrow{kpi}{balance}{mass\_closure\_rel}{4.0063377475e-10}
\goldrow{kpi}{balance}{mass\_kg\_final}{10701.9131449}
\goldrow{kpi}{balance}{mass\_kg\_in\_cum}{5443104}
\goldrow{kpi}{balance}{mass\_kg\_initial}{4989.512}
\goldrow{kpi}{balance}{mass\_kg\_out\_cum}{5437391.60103}
\goldrow{kpi}{balance}{mass\_residual\_kg}{0.00218069130187}
\goldrow{kpi}{balance}{material\_available}{1}
\goldrow{kpi}{plant}{F\_pP\_klbh\_final}{3.90161524936}
\goldrow{kpi}{plant}{F\_wG\_klbh\_final}{1.22338475064}
\goldrow{kpi}{plant}{J\_combined\_klb\_final}{1071.35532885}
\goldrow{kpi}{plant}{T\_final}{322.2222}
\goldrow{kpi}{plant}{t\_end}{1440000}
\goldrow{kpi}{plant}{waste\_klb\_final}{481.451273663}
\goldmore{14}
\end{tutgolden}

\clearpage
\tutentry{ctrl13\_williams\_otto\_step}{tutorials/ctrl/ctrl13\_williams\_otto\_step}
\begin{tutscheme}
\begin{tikzpicture}[>={Stealth[length=2mm]}, every node/.style={font=\scriptsize}]
  \node[draw, rounded corners=2pt, align=center, minimum width=1.9cm, minimum height=0.85cm, fill=black!4] (plant) at (0.0,-0.0) {\textbf{plant}\\[-1pt]{\tiny williamsOttoPlant}};
\end{tikzpicture}
\end{tutscheme}
\tutwhat{Williams-Otto F\_fB step 20$\to$21 klb/h at t$\ast$=100 h (paper Fig. 2): F\_pP climbs \textasciitilde{}3.90 $\to$ \textasciitilde{}4.45 klb/h}
\begin{tutnarrative}
One unit, because the benchmark IS one equation set: the Williams-Otto reactor + decanter + column + splitter fold into six ODEs (the recycle term (1-eta)mu, the removed G, the 0.1 F\_tE column retention are all inside them -- see the unit's own header). The five controls sit in `operation`, declared SI: with the MW 453.592 convention (1 kmol == 1 klb) every numeral below IS the paper's value in klb/h.

u$\ast$ of eq. (3.20) plus the paper's FIRST simulation experiment (section 3.2, Fig. 2): a Schedule steps F\_fB from 20 to 21 klb/h at t$\ast$ = 100 h. The MV crosses the boundary in SI (kmol/s), so the step values below are 20/3600 and 21/3600 -- the paper's numbers over the one declared conversion.

ONE DECLARED FIGURE (system/postDict, C487): `response`, "how do the product and waste rates answer a 5 \% step in the feed of B?" -- both rates (klb/h) against hours, the stepped feed on the right-hand axis.

\end{tutnarrative}
\begin{tutgolden}
\goldrow{kpi}{FeedStep}{MV\_final}{0.0058333333}
\goldrow{kpi}{balance}{energy\_balance\_available}{0}
\goldrow{kpi}{balance}{mass\_closure\_rel}{1.45035580786e-09}
\goldrow{kpi}{balance}{mass\_kg\_final}{19794.220051}
\goldrow{kpi}{balance}{mass\_kg\_in\_cum}{2766911.20181}
\goldrow{kpi}{balance}{mass\_kg\_initial}{4989.512}
\goldrow{kpi}{balance}{mass\_kg\_out\_cum}{2752106.49778}
\goldrow{kpi}{balance}{mass\_residual\_kg}{0.00401300573139}
\goldrow{kpi}{balance}{material\_available}{1}
\goldrow{kpi}{plant}{F\_pP\_klbh\_final}{4.44568758172}
\goldrow{kpi}{plant}{F\_wG\_klbh\_final}{1.90791543668}
\goldrow{kpi}{plant}{J\_combined\_klb\_final}{529.364471374}
\goldrow{kpi}{plant}{T\_final}{322.2222}
\goldrow{kpi}{plant}{t\_end}{720000}
\goldmore{9}
\end{tutgolden}

\clearpage
\tutentry{ctrl14\_williams\_otto\_pi}{tutorials/ctrl/ctrl14\_williams\_otto\_pi}
\begin{tutscheme}
\begin{tikzpicture}[>={Stealth[length=2mm]}, every node/.style={font=\scriptsize}]
  \node[draw, rounded corners=2pt, align=center, minimum width=1.9cm, minimum height=0.85cm, fill=black!4] (plant) at (0.0,-0.0) {\textbf{plant}\\[-1pt]{\tiny williamsOttoPlant}};
\end{tikzpicture}
\end{tutscheme}
\tutwhat{Williams-Otto PI channel (F\_fB, F\_pP), paper eq. 4.6 with the section-4.1 tunings: F\_pP tracks 3.90 $\to$ 4.5 klb/h with the Fig. 4 overshoot}
\begin{tutnarrative}
One unit, because the benchmark IS one equation set: the Williams-Otto reactor + decanter + column + splitter fold into six ODEs (the recycle term (1-eta)mu, the removed G, the 0.1 F\_tE column retention are all inside them -- see the unit's own header). The five controls sit in `operation`, declared SI: with the MW 453.592 convention (1 kmol == 1 klb) every numeral below IS the paper's value in klb/h.

The paper's control channel (F\_fB, F\_pP), eq. (4.6), with the Skogestad-IMC gains of section 4.1: K3prop = 0.069 (dimensionless -- klb/h per klb/h survives the SI conversion unchanged) and K3int = 1.282 per HOUR, declared per second (1.282/3600 = 3.5611111e-4). Bias = F$\ast$\_fB (bumpless), setpoint F\_pP = 4.5 klb/h = 1.25e-3 kmol/s, output bounded to the paper's 0 <= F\_fB <= 56 klb/h.

ONE DECLARED FIGURE (system/postDict, C487): `response`, "does the PI loop bring the product rate to its set point, and how far does the feed of B have to move?" -- PV against the set point, the manipulated feed on the right-hand axis, against hours.

\end{tutnarrative}
\begin{tutgolden}
\goldrow{kpi}{PC1}{MV\_final}{0.00586183454679}
\goldrow{kpi}{PC1}{PV\_final}{0.00124999896025}
\goldrow{kpi}{PC1}{SP\_final}{0.00125}
\goldrow{kpi}{balance}{energy\_balance\_available}{0}
\goldrow{kpi}{balance}{mass\_closure\_rel}{2.05158152402e-09}
\goldrow{kpi}{balance}{mass\_kg\_final}{21425.7971042}
\goldrow{kpi}{balance}{mass\_kg\_in\_cum}{1418738.01455}
\goldrow{kpi}{balance}{mass\_kg\_initial}{10701.9131449}
\goldrow{kpi}{balance}{mass\_kg\_out\_cum}{1408014.1335}
\goldrow{kpi}{balance}{mass\_residual\_kg}{0.00291065669808}
\goldrow{kpi}{balance}{material\_available}{1}
\goldrow{kpi}{plant}{F\_pP\_klbh\_final}{4.49999626428}
\goldrow{kpi}{plant}{F\_wG\_klbh\_final}{2.00328828826}
\goldrow{kpi}{plant}{J\_combined\_klb\_final}{255.611908728}
\goldmore{11}
\end{tutgolden}

\clearpage
\tutentry{ctrl16\_williams\_otto\_optimal}{tutorials/ctrl/ctrl16\_williams\_otto\_optimal}
\begin{tutscheme}
\begin{tikzpicture}[>={Stealth[length=2mm]}, every node/.style={font=\scriptsize}]
  \node[draw, rounded corners=2pt, align=center, minimum width=1.9cm, minimum height=0.85cm, fill=black!4] (plant) at (0.0,-0.0) {\textbf{plant}\\[-1pt]{\tiny williamsOttoPlant}};
\end{tikzpicture}
\end{tutscheme}
\tutwhat{Williams-Otto section-5.3 combined optimisation: Nelder-Mead over two K=5 Schedule profiles, J measured against the published collocation optimum 551.8}
\begin{tutnarrative}
One unit, because the benchmark IS one equation set: the Williams-Otto reactor + decanter + column + splitter fold into six ODEs (the recycle term (1-eta)mu, the removed G, the 0.1 F\_tE column retention are all inside them -- see the unit's own header). The five controls sit in `operation`, declared SI: with the MW 453.592 convention (1 kmol == 1 klb) every numeral below IS the paper's value in klb/h.

The section-5.3 combined optimisation: system/outerDict shapes the two profiles below to maximise J = yield - waste (alpha = beta = 1, eq. 5.7), read straight off the plant's J\_combined\_klb\_final KPI.

\end{tutnarrative}
\begin{tutgolden}
\goldrow{kpi}{FfBprof}{MV\_final}{0.00463988515428}
\goldrow{kpi}{Tprof}{MV\_final}{300.373635533}
\goldrow{kpi}{balance}{energy\_balance\_available}{0}
\goldrow{kpi}{balance}{mass\_closure\_rel}{1.1440186567e-07}
\goldrow{kpi}{balance}{mass\_kg\_final}{723583.798797}
\goldrow{kpi}{balance}{mass\_kg\_in\_cum}{2128628.70559}
\goldrow{kpi}{balance}{mass\_kg\_initial}{10701.9131449}
\goldrow{kpi}{balance}{mass\_kg\_out\_cum}{1415747.06346}
\goldrow{kpi}{balance}{mass\_residual\_kg}{0.243519095238}
\goldrow{kpi}{balance}{material\_available}{1}
\goldrow{kpi}{plant}{F\_pP\_klbh\_final}{7.05702728738}
\goldrow{kpi}{plant}{F\_wG\_klbh\_final}{1.47544148411}
\goldrow{kpi}{plant}{J\_combined\_klb\_final}{546.81259418}
\goldrow{kpi}{plant}{T\_final}{300.373635533}
\goldmore{33}
\end{tutgolden}

\clearpage
\tutentry{ctrl17\_signal\_shapes}{tutorials/ctrl/ctrl17\_signal\_shapes}
\begin{tutscheme}
\begin{tikzpicture}[>={Stealth[length=2mm]}, every node/.style={font=\scriptsize}]
  \node[draw, rounded corners=2pt, align=center, minimum width=1.9cm, minimum height=0.85cm, fill=black!4] (reactor) at (0.0,-0.0) {\textbf{reactor}\\[-1pt]{\tiny dynamicCSTR}};
  \draw[->] ($(reactor.west)+(-1.05,0.00)$) -- ($(reactor.west)+(0,0.00)$) node[at start, above, font=\tiny, anchor=south west] {feed};
  \draw[->] ($(reactor.east)+(0,0.00)$) -- ($(reactor.east)+(1.05,0.00)$) node[at end, above, font=\tiny, anchor=south east] {product};
\end{tikzpicture}
\end{tutscheme}
\tutwhat{WITNESS: ramp + pulse + staircase signals on three feed handles, with the temperature PID rejecting all three at once; the three registered shapes no case selected.}
\begin{tutnarrative}
WITNESS: the ramp, pulse and staircase SIGNALS, each on its own feed handle of the ctrl09 CSTR, with the ctrl09 temperature PID still in the loop rejecting all three at once -- compA ramps +28.6 \% over 1000 s, compB takes a 300 s rectangular pulse, the feed temperature walks a three-step staircase. These three registered signal shapes were selected by no case until 2026-08-23 (invariant I8). Self-recorded.

TWO DECLARED FIGURES (system/postDict, C487): `loop`, "does the temperature loop hold the reactor at 350 K while three different disturbances arrive at once?" (from 400 s, tight axis 349.7 to 350.2 K), and `feeds`, the three disturbances themselves.

\end{tutnarrative}
\begin{tutgolden}
\goldrow{kpi}{FeedPulseB}{MV\_final}{1.5e-05}
\goldrow{kpi}{FeedPulseB}{PV\_final}{0}
\goldrow{kpi}{FeedPulseB}{SP\_final}{1.5e-05}
\goldrow{kpi}{FeedRampA}{MV\_final}{4.5e-05}
\goldrow{kpi}{FeedRampA}{PV\_final}{0}
\goldrow{kpi}{FeedRampA}{SP\_final}{4.5e-05}
\goldrow{kpi}{FeedStairT}{MV\_final}{348}
\goldrow{kpi}{FeedStairT}{PV\_final}{0}
\goldrow{kpi}{FeedStairT}{SP\_final}{348}
\goldrow{kpi}{TC1}{MV\_final}{347.378942902}
\goldrow{kpi}{TC1}{PV\_final}{349.993146257}
\goldrow{kpi}{TC1}{SP\_final}{350}
\goldrow{kpi}{balance}{energy\_balance\_available}{0}
\goldrow{kpi}{balance}{mass\_closure\_rel}{2.98582911041e-14}
\goldmore{13}
\end{tutgolden}

\clearpage
\tutentry{ctrl17\_step\_reaction\_curve}{tutorials/ctrl/ctrl17\_step\_reaction\_curve}
\begin{tutscheme}
\begin{tikzpicture}[>={Stealth[length=2mm]}, every node/.style={font=\scriptsize}]
  \node[draw, rounded corners=2pt, align=center, minimum width=1.9cm, minimum height=0.85cm, fill=black!4] (reactor) at (0.0,-0.0) {\textbf{reactor}\\[-1pt]{\tiny dynamicCSTR}};
  \draw[->] ($(reactor.west)+(-1.05,0.00)$) -- ($(reactor.west)+(0,0.00)$) node[at start, above, font=\tiny, anchor=south west] {feed};
  \draw[->] ($(reactor.east)+(0,0.00)$) -- ($(reactor.east)+(1.05,0.00)$) node[at end, above, font=\tiny, anchor=south east] {product};
\end{tikzpicture}
\end{tutscheme}
\tutwhat{Open-loop process reaction curve: the jacket steps +5, -5, +20 and -20 K from one settled base, and the four responses are the identification record for the FOPDT fit in system/propsDict}
\begin{tutnarrative}
THE IDENTIFICATION, run over the trajectory the choupoCtrl half of this case wrote. Run the experiment first:

runCase tutorials/ctrl/ctrl17\_step\_reaction\_curve choupoProps tutorials/ctrl/ctrl17\_step\_reaction\_curve

Each operation DECLARES that its trajectory is the process run's (`writtenBy \{ trajectory process; \}`), so the second command run alone, on a fresh copy, refuses by name and says to run the first.

TWO OPERATIONS OVER ONE EXPERIMENT, and the pair is the lesson. Both read the SAME four steps of the SAME jacket temperature; they differ only in which response they call the controlled variable.

[0] the COMPOSITION loop (jacket $\to$ product holdup). Two lags in series -- the jacket heats the liquid, and the warmer liquid reacts faster -- so the response is S-shaped and an FOPDT surrogate finds a real apparent dead time. All three tuning rules apply.

[1] the TEMPERATURE loop (jacket $\to$ reactor temperature). One dominant lag and no transport delay, so the identified theta comes out smaller than the trajectory's own sampling interval. The engine says so, and REFUSES every Ziegler-Nichols and Cohen-Coon setting by name: both rules put theta in a denominator, and a controller gain computed from an unresolved number is a number with nothing behind it. IMC survives, because its algebra never divides by theta -- which is exactly why it is the rule of choice for a lag-dominant loop.

Nothing was tuned to make that happen. It is one plant, one experiment, and two honest answers, and the difference between them is the whole reason the reaction-curve method needs a dead time to work at all.

\end{tutnarrative}
\begin{tutgolden}
\goldrow{kpi}{JacketStep}{MV\_final}{320}
\goldrow{kpi}{reactor}{T\_final}{320.759149686}
\goldrow{kpi}{reactor}{n\_compA\_final}{0.0117959393424}
\goldrow{kpi}{reactor}{n\_compB\_final}{0.000204060657576}
\goldrow{kpi}{reactor}{F\_in\_final}{5e-05}
\goldrow{kpi}{reactor}{t\_end}{15000}
\goldrow{kpi}{balance}{energy\_balance\_available}{0}
\goldrow{kpi}{balance}{mass\_closure\_rel}{9.91527847481e-14}
\goldrow{kpi}{balance}{mass\_kg\_final}{0.36}
\goldrow{kpi}{balance}{mass\_kg\_in\_cum}{22.5}
\goldrow{kpi}{balance}{mass\_kg\_initial}{0.36}
\goldrow{kpi}{balance}{mass\_kg\_out\_cum}{22.5}
\goldrow{kpi}{balance}{mass\_residual\_kg}{2.23093765683e-12}
\goldrow{kpi}{balance}{material\_available}{1}
\goldmore{2}
\end{tutgolden}

\clearpage
\tutentry{ctrl18\_rtd\_pulse\_cstr}{tutorials/ctrl/ctrl18\_rtd\_pulse\_cstr}
\begin{tutscheme}
\begin{tikzpicture}[>={Stealth[length=2mm]}, every node/.style={font=\scriptsize}]
  \node[draw, rounded corners=2pt, align=center, minimum width=1.9cm, minimum height=0.85cm, fill=black!4] (reactor) at (0.0,-0.0) {\textbf{reactor}\\[-1pt]{\tiny dynamicCSTR}};
  \draw[->] ($(reactor.west)+(-1.05,0.00)$) -- ($(reactor.west)+(0,0.00)$) node[at start, above, font=\tiny, anchor=south west] {feed};
  \draw[->] ($(reactor.east)+(0,0.00)$) -- ($(reactor.east)+(1.05,0.00)$) node[at end, above, font=\tiny, anchor=south east] {product};
\end{tikzpicture}
\end{tutscheme}
\tutwhat{RTD by pulse tracer on a pure mixing tank: tracer substitution at constant flow, tau = N/F = 240 s; engine-measured moments against Levenspiel closed-form anchors.}
\begin{tutnarrative}
RESIDENCE-TIME DISTRIBUTION by pulse tracer -- the engine's first RTD tool (2026-08-23). A pure mixing tank (the ctrl09 vessel, reaction and controller stripped) holds N = 0.012 kmol and is fed F = 5e-5 kmol/s, so tau = N/F = 240 s exactly. For 30 s starting at t = 100 s the feed SWAPS 1e-5 kmol/s of carrier for tracer -- compB pulses up by exactly what compA pulses down -- so the total flow never moves and the vessel is rigorously linear and time-invariant in the tracer.

THE ANCHORS ARE CLOSED FORMS (Levenspiel, Chemical Reaction Engineering, ch. 13-14 -- moment additivity of a linear vessel, exact for ANY inlet shape): tbar\_out = tbar\_in + tau = (100 + 15) + 240 = 355 s sigma2\_out = sigma2\_in + tau\textasciicircum{}2 = 30\textasciicircum{}2/12 + 240\textasciicircum{}2 = 57675 s\textasciicircum{}2 area = amplitude x width = 1e-5 $\ast$ 30 = 3.0e-4 kmol The engine measures the outlet response moments (rtd KPIs, trapezoid on the accepted-step grid, reports/rtd/E.csv); the arithmetic above is the case's, stated here. Deviations from the anchors are integration truncation (run = 2600 s \textasciitilde{} 10.4 tau past the pulse; e\textasciicircum{}-10.4 \textasciitilde{} 3e-5) plus the finite step -- both stated, neither hidden.

TWO DECLARED FIGURES (system/postDict, C487): `pulse`, "when does the tracer pulse enter the tank, and when has it left?", and `washout`, from 130 s the tracer in the tank falling by a factor e every 240 s.

\end{tutnarrative}
\begin{tutgolden}
\goldrow{kpi}{CarrierDip}{MV\_final}{5e-05}
\goldrow{kpi}{CarrierDip}{PV\_final}{0}
\goldrow{kpi}{CarrierDip}{SP\_final}{5e-05}
\goldrow{kpi}{TracerPulse}{MV\_final}{0}
\goldrow{kpi}{TracerPulse}{PV\_final}{0}
\goldrow{kpi}{TracerPulse}{SP\_final}{0}
\goldrow{kpi}{balance}{energy\_balance\_available}{0}
\goldrow{kpi}{balance}{mass\_closure\_rel}{4.73695157173e-14}
\goldrow{kpi}{balance}{mass\_kg\_final}{0.36}
\goldrow{kpi}{balance}{mass\_kg\_in\_cum}{3.9}
\goldrow{kpi}{balance}{mass\_kg\_initial}{0.36}
\goldrow{kpi}{balance}{mass\_kg\_out\_cum}{3.9}
\goldrow{kpi}{balance}{mass\_residual\_kg}{1.84741111298e-13}
\goldrow{kpi}{balance}{material\_available}{1}
\goldmore{15}
\end{tutgolden}

\clearpage
\tutentry{ctrl19\_freq\_response\_cstr}{tutorials/ctrl/ctrl19\_freq\_response\_cstr}
\begin{tutscheme}
\begin{tikzpicture}[>={Stealth[length=2mm]}, every node/.style={font=\scriptsize}]
  \node[draw, rounded corners=2pt, align=center, minimum width=1.9cm, minimum height=0.85cm, fill=black!4] (reactor) at (0.0,-0.0) {\textbf{reactor}\\[-1pt]{\tiny dynamicCSTR}};
  \draw[->] ($(reactor.west)+(-1.05,0.00)$) -- ($(reactor.west)+(0,0.00)$) node[at start, above, font=\tiny, anchor=south west] {feed};
  \draw[->] ($(reactor.east)+(0,0.00)$) -- ($(reactor.east)+(1.05,0.00)$) node[at end, above, font=\tiny, anchor=south east] {product};
\end{tikzpicture}
\end{tutscheme}
\tutwhat{Frequency response of the LTI mixing tank at omega$\ast$tau = 1: engine-fitted amplitude and phase against the exact first-order closed forms (AR = 1/sqrt(2), phase = -45 deg).}
\begin{tutnarrative}
FREQUENCY RESPONSE of the same LTI mixing tank (tau = N/F = 240 s): the tracer feed oscillates sinusoidally at omega$\ast$tau = 1 -- again as a SUBSTITUTION against the carrier, so total flow never moves and the vessel is rigorously first order in the tracer. The engine fits a + b sin + c cos over 4 cycles after discarding 5 (transient e\textasciicircum{}-t/tau through 5 periods = e\textasciicircum{}-31 -- gone), and reports amplitude, phase and the unexplained-variance residual.

THE ANCHORS ARE THE FIRST-ORDER CLOSED FORMS (any process-control text; at omega$\ast$tau = 1 exactly): AR = 1/sqrt(1 + (omega tau)\textasciicircum{}2) = 1/sqrt(2) $\to$ out amplitude = 2.0e-6 $\ast$ 0.7071068 = 1.4142136e-6 kmol/s phase = -atan(omega tau) = -45 deg = -0.7853982 rad (the drive is sin(w(t-tStart)) with tStart 0, so the fitted phase IS the lag). The measured amplitude lands within 1e-6 relative of the closed form and the fit residual is \textasciitilde{}7e-15 (the vessel really is LTI); the phase carries a bias of EXACTLY omega$\ast$dt/2 = 0.00208 rad = 0.12 deg -- the half-sample lag of reading accepted states on a 1 s grid, the measurement system's own phase contribution, worth teaching in itself. The anchor band covers it and this header names it.

ONE DECLARED FIGURE (system/postDict, C487): `response`, "how much does the tank damp and delay a sinusoidal tracer feed?" -- the tracer in the tank against the tracer feed.

\end{tutnarrative}
\begin{tutgolden}
\goldrow{kpi}{CarrierMirror}{MV\_final}{4.41011922784e-05}
\goldrow{kpi}{CarrierMirror}{PV\_final}{0}
\goldrow{kpi}{CarrierMirror}{SP\_final}{4.41011922784e-05}
\goldrow{kpi}{TracerSine}{MV\_final}{5.89880772158e-06}
\goldrow{kpi}{TracerSine}{PV\_final}{0}
\goldrow{kpi}{TracerSine}{SP\_final}{5.89880772158e-06}
\goldrow{kpi}{balance}{energy\_balance\_available}{0}
\goldrow{kpi}{balance}{mass\_closure\_rel}{2.91033739366e-15}
\goldrow{kpi}{balance}{mass\_kg\_final}{0.36}
\goldrow{kpi}{balance}{mass\_kg\_in\_cum}{20.8095}
\goldrow{kpi}{balance}{mass\_kg\_initial}{0.36}
\goldrow{kpi}{balance}{mass\_kg\_out\_cum}{20.8095}
\goldrow{kpi}{balance}{mass\_residual\_kg}{-6.05626659933e-14}
\goldrow{kpi}{balance}{material\_available}{1}
\goldmore{14}
\end{tutgolden}

\clearpage
\tutentry{ctrl20\_bode\_cstr}{tutorials/ctrl/ctrl20\_bode\_cstr}
\begin{tutscheme}
\begin{tikzpicture}[>={Stealth[length=2mm]}, every node/.style={font=\scriptsize}]
  \node[draw, rounded corners=2pt, align=center, minimum width=1.9cm, minimum height=0.85cm, fill=black!4] (reactor) at (0.0,-0.0) {\textbf{reactor}\\[-1pt]{\tiny dynamicCSTR}};
  \draw[->] ($(reactor.west)+(-1.05,0.00)$) -- ($(reactor.west)+(0,0.00)$) node[at start, above, font=\tiny, anchor=south west] {feed};
  \draw[->] ($(reactor.east)+(0,0.00)$) -- ($(reactor.east)+(1.05,0.00)$) node[at end, above, font=\tiny, anchor=south east] {product};
\end{tikzpicture}
\end{tutscheme}
\tutwhat{Bode table of the LTI mixing tank: outer frequency sweep (one datum, lockstep mirror) over the engine frequency-response fit; anchors are the exact first-order closed forms at each point.}
\begin{tutnarrative}
THE BODE TABLE of the LTI mixing tank (tau = 240 s): an outer sweep drives the sinusoidal tracer substitution at four frequencies, omega$\ast$tau = 0.5, 2.0, 3.5, 5.0 (linear grid in f), and the engine's frequencyResponse fit measures amplitude and phase at each -- the sweep CSV IS the Bode table. The measurement block references the drive, so each pass's fit window adapts to that pass's own period, and the sweep varies ONE datum (the drive frequency, mirrored to the carrier by linkedTargets -- the same value, lockstep, never a second independent number).

ANCHORS: the exact first-order closed forms at each grid point, AR = 1/sqrt(1+(wt)\textasciicircum{}2), phase = -atan(wt): wt=0.5: |out| = 2e-6$\ast$0.894427 = 1.788854e-6, phi = -0.463648 rad wt=2.0: |out| = 2e-6$\ast$0.447214 = 8.944272e-7, phi = -1.107149 rad wt=3.5: |out| = 2e-6$\ast$0.274721 = 5.494423e-7, phi = -1.292497 rad wt=5.0: |out| = 2e-6$\ast$0.196116 = 3.922323e-7, phi = -1.373401 rad Phase bands cover the omega$\ast$dt/2 half-sample bias, which grows with omega (0.10 to 1.0 mrad here) -- stated, not hidden.

\end{tutnarrative}
\begin{tutgolden}
\goldrow{csv}{sweep\_results.csv}{0:frequencyResponse.out\_amplitude\_kmol\_s}{1.788854e-6}
\goldrow{csv}{sweep\_results.csv}{1:frequencyResponse.out\_amplitude\_kmol\_s}{8.944272e-7}
\goldrow{csv}{sweep\_results.csv}{2:frequencyResponse.out\_amplitude\_kmol\_s}{5.494423e-7}
\goldrow{csv}{sweep\_results.csv}{3:frequencyResponse.out\_amplitude\_kmol\_s}{3.922323e-7}
\goldrow{csv}{sweep\_results.csv}{0:frequencyResponse.out\_phase\_rad}{-0.463648}
\goldrow{csv}{sweep\_results.csv}{1:frequencyResponse.out\_phase\_rad}{-1.107149}
\goldrow{csv}{sweep\_results.csv}{2:frequencyResponse.out\_phase\_rad}{-1.292497}
\goldrow{csv}{sweep\_results.csv}{3:frequencyResponse.out\_phase\_rad}{-1.373401}
\goldrow{kpi}{CarrierMirror}{MV\_final}{4.79049069948e-05}
\goldrow{kpi}{CarrierMirror}{PV\_final}{0}
\goldrow{kpi}{CarrierMirror}{SP\_final}{4.79049069948e-05}
\goldrow{kpi}{TracerSine}{MV\_final}{2.09509300524e-06}
\goldrow{kpi}{TracerSine}{PV\_final}{0}
\goldrow{kpi}{TracerSine}{SP\_final}{2.09509300524e-06}
\goldmore{20}
\end{tutgolden}

\clearpage
\tutentry{ctrl21\_tank\_level\_control}{tutorials/ctrl/ctrl21\_tank\_level\_control}
\begin{tutscheme}
\begin{tikzpicture}[>={Stealth[length=2mm]}, every node/.style={font=\scriptsize}]
  \node[draw, rounded corners=2pt, align=center, minimum width=1.9cm, minimum height=0.85cm, fill=black!4] (tank) at (0.0,-0.0) {\textbf{tank}\\[-1pt]{\tiny dynamicTank}};
  \draw[->] ($(tank.west)+(-1.05,0.00)$) -- ($(tank.west)+(0,0.00)$) node[at start, above, font=\tiny, anchor=south west] {feed};
  \draw[->] ($(tank.east)+(0,0.00)$) -- ($(tank.east)+(1.05,0.00)$) node[at end, above, font=\tiny, anchor=south east] {product};
\end{tikzpicture}
\end{tutscheme}
\tutwhat{Water tank with a manipulated outflow and a PI level controller: the inflow steps +50 \% at 600 s, the integrating level is brought back to the 2.0 m setpoint (P-only would settle 0.05 m off)}
\begin{tutnarrative}
the most common loop of a plant: INVENTORY.

The same water tank as unsteady05 (D = 2 m, area 3.1416 m2), but its outflow is a MANIPULATED FLOW (a flow-controlled pump: `mode flow`, mv F\_out), so the level INTEGRATES the difference of the flows,

A rho/MW dh/dt = F\_in - F\_out plant gain a = MW/(rho A) = 5.7e-3 m/s per kmol/s (rho = 1001.55 kg/m3)

and nothing brings it back: an integrating process. A level controller closes the loop on F\_out. The inflow steps at t = 600 s from 0.20 to 0.30 kmol/s (a Signal on the tank's inlet face).

THE SIGN. The controller is direct-acting on error = setpoint - level, and a level ABOVE the setpoint must RAISE the outflow, so the gains are NEGATIVE (u = bias + Kp (SP - h) + Ki Int(SP - h) with Kp < 0). Bias 0.20 kmol/s is the outflow that balances the old inflow.

WHAT THE TWO CONTROLLERS DO (measured; check\_dynamic\_tank re-runs both):

PI (Kp = -2, Ki = -0.01): the integral term removes the offset; the level returns to the 2.0 m setpoint (zero offset), after a peak excursion of 0.031 m (max |h - 2| in trajectory.csv; the end offset is < 1e-7 m). P-only (Ki = 0, edit one line): the outflow must rise by 0.10 kmol/s and a proportional-only controller can only do that with an error, F\_out = bias + Kp (SP - h) $\to$ h - SP = (F\_in,new - bias)/|Kp| = 0.10 / 2 = 0.050 m (offset: the level settles at 2.05 m).

The closed loop is second order, wn\textasciicircum{}2 = a |Ki| , 2 zeta wn = a |Kp|: wn = 7.6e-3 1/s, zeta = 0.76. The inflow is a disturbance the loop must reject; the tank's own buffering (volume) is what the controller trades against the downstream flow's smoothness -- loose level control is a feature, not a defect, when the outflow feeds another unit.

ENERGY. H on the elements datum, T a readout; the ledger claims mass and energy and closes (check\_dynamic\_tank).

ONE DECLARED FIGURE (system/postDict, C487): `level`, "does the level controller bring the tank back to its set point after the inflow steps up?" -- the level against its set point, the inflow and the manipulated outflow on the right-hand axis.

\end{tutnarrative}
\tutreadme{A water tank (dynamicTank, D = 2 m, 3.14 m$^2$) whose outflow is a manipulated flow (mode flow, mv F\_out, a flow-controlled pump). Because nothing in the outflow depends on the level, the level integrates the difference of the flows and never comes back by itself. A PI level controller on F\_out closes the loop. The inflow steps from 0.20 to 0.30 kmol/s at t = 600 s; 3000 s run, 2 s steps.}
\tutreadme{Run it: runCase tutorials/ctrl/ctrl21\_tank\_level\_control (it runs under choupoCtrl). Its sibling unsteady05\_tank\_valve\_drain drains the same tank through a valve instead and needs no controller.}
\tutreadme{1. The PI holds the level. Plot level in trajectory.csv: it stays at the 2.0 m setpoint, rises by at most 0.031 m after the inflow step and returns with zero offset (< 1e-7 m at the end). The integral term is what removes the offset. 2. P-only settles off the setpoint. Set Ki 0.0; in the controller's gains \{\} and run again: the level ends at 2.05 m. A proportional controller needs an error to hold a different outflow, and the offset is (F\_in,new - bias)/ = (0.30 - 0.20)/2 = 0.050 m. 3. The gains are negative. The controller is direct-acting on error = setpoint - level; a level ABOVE the setpoint must RAISE the outflow, so the loop is reverse-acting and Kp, Ki are negative. With a positive sign the level runs away until the tank overflows, and the engine refuses that by name (OVERFLOWED). 4. The ledger claims mass and energy and closes (\textasciitilde{}1e-14): the tank stores its enthalpy on the elements datum, T is a readout. [balance] lines print the check.}
\tutreadme{The closed loop is second order (?n $\approx$ 7.6e-3 1/s, ? $\approx$ 0.76). The tank's buffering is what a loose level loop trades against the smoothness of the downstream flow. check\_dynamic\_tank re-runs both the PI and the P-only variants and holds these numbers.}
\begin{tutgolden}
\goldrow{kpi}{FeedStep}{MV\_final}{0.3}
\goldrow{kpi}{FeedStep}{PV\_final}{0}
\goldrow{kpi}{FeedStep}{SP\_final}{0.3}
\goldrow{kpi}{LC1}{MV\_final}{0.299999850013}
\goldrow{kpi}{LC1}{PV\_final}{1.99999993576}
\goldrow{kpi}{LC1}{SP\_final}{2}
\goldrow{kpi}{balance}{element\_H\_closure\_rel}{2.31433916649e-14}
\goldrow{kpi}{balance}{element\_O\_closure\_rel}{2.31433916649e-14}
\goldrow{kpi}{balance}{element\_worst\_closure\_rel}{2.31433916649e-14}
\goldrow{kpi}{balance}{energy\_balance\_available}{1}
\goldrow{kpi}{balance}{energy\_closure\_rel}{3.79805900307e-14}
\goldrow{kpi}{balance}{energy\_functional\_H}{1}
\goldrow{kpi}{balance}{energy\_kJ\_final}{-99842109.8295}
\goldrow{kpi}{balance}{energy\_kJ\_heat\_cum}{0}
\goldmore{20}
\end{tutgolden}

\clearpage
\tutentry{ctrl22\_reactor\_level\_and\_temperature}{tutorials/ctrl/ctrl22\_reactor\_level\_and\_temperature}
\begin{tutscheme}
\begin{tikzpicture}[>={Stealth[length=2mm]}, every node/.style={font=\scriptsize}]
  \node[draw, rounded corners=2pt, align=center, minimum width=1.9cm, minimum height=0.85cm, fill=black!4] (reactor) at (0.0,-0.0) {\textbf{reactor}\\[-1pt]{\tiny dynamicCSTR}};
  \draw[->] ($(reactor.west)+(-1.05,0.00)$) -- ($(reactor.west)+(0,0.00)$) node[at start, above, font=\tiny, anchor=south west] {feed};
  \draw[->] ($(reactor.east)+(0,0.00)$) -- ($(reactor.east)+(1.05,0.00)$) node[at end, above, font=\tiny, anchor=south east] {product};
\end{tikzpicture}
\end{tutscheme}
\tutwhat{Reactor with a level: a PI level loop on the outflow and a PI temperature loop on the jacket act on ONE unit while the inflow steps +50 \% -- the level returns to its setpoint, the temperature holds, the jacket area follows the liquid}
\begin{tutnarrative}
A jacketed reactor with a LEVEL (dynamicCSTR, `volumeMode variable`, C205 slice 2): isobutanol $\to$ 2-butanol, first order, exothermic, fed at 300 K and run at 340 K. Its outflow is a MANIPULATED FLOW (as in ctrl21), so the level INTEGRATES the flow imbalance and needs a controller; its temperature is held by the jacket (as in ctrl01). The two loops act on different manipulated variables of the same vessel:

LC1 PI cv level $\to$ mv F\_out reverse acting (Kp, Ki < 0) TC1 PI cv T $\to$ mv T\_jacket direct acting on the error T\_SP - T

and they are COUPLED through the vessel, which is the lesson:

the level sets the volume, so the jacket's wetted area -- the declared UA is the value with the vessel FULL, UA(h) = UA A\_w(h)/A\_w(H) with A\_w = pi D h + pi D\textasciicircum{}2/4 -- and the reaction heat r V both follow it; the outflow leaves at the vessel's temperature, so it carries enthalpy away; the canonical energy route (announced at start-up) lets that term cancel exactly, and the ledger closes for the whole run.

The inflow steps from 1.0e-3 to 1.5e-3 kmol/s at t = 1200 s (a Signal on the inlet face). Level plant gain a = v/A = 0.47 m per kmol (liquid molar volume 0.093 m3/kmol over 0.196 m2); with wn = 0.01 1/s, zeta = 0.8 the PI gains are Kp = -2 zeta wn / a = -0.034, Ki = -wn\textasciicircum{}2 / a = -2.1e-4 kmol/s per m(s). Bias 1.0e-3 kmol/s is the outflow that balances the old inflow.

The cold feed lowers the temperature as the inflow rises; the jacket absorbs it. Read trajectory.csv: `level`, `T`, `F\_out`, `volume`, and the controllers' MV columns.

TWO DECLARED FIGURES (system/postDict, C487): `level`, "does the level loop hold the reactor's liquid when the inflow steps up?", and `temperature`, "does the temperature loop hold the reactor at 340 K through the start-up and the step?" (the jacket temperature on the right-hand axis).

\end{tutnarrative}
\tutreadme{A jacketed reactor with a level (dynamicCSTR with volumeMode variable, D = 0.5 m, H = 1.0 m): isobutanol $\to$ 2-butanol, first order, exothermic, fed at 300 K and run at 340 K. Two PI controllers act on DIFFERENT manipulated variables of the SAME vessel:}
\tutreadme{$\ast$ LC1 --- level on the outflow F\_out, reverse acting (Kp = -0.034, Ki = -2.1e-4); $\ast$ TC1 --- temperature on the jacket T\_jacket (Kp = 5, Ki = 0.02).}
\tutreadme{The inflow steps from 1.0e-3 to 1.5e-3 kmol/s at t = 1200 s; 3600 s run, 2 s steps. Run it: runCase tutorials/ctrl/ctrl22\_reactor\_level\_and\_temperature (it runs under choupoCtrl). Its sibling unsteady06\_reactor\_fill\_and\_react fills the same reactor from a heel with no controller.}
\tutreadme{1. The level loop rejects the inflow step. Plot reactor.level: the peak excursion is 0.010 m (the P-only bound would be dF/ = 0.015 m) and the level is back on 0.50 m at the end (< 2e-6 m). 2. The temperature loop holds 340 K against a colder, larger feed: T dips about 2 K after the step and returns; the jacket temperature (TC1.MV) rises from about 367 K to about 387 K. 3. The loops are coupled through the vessel. The declared UA = 300 W/K is the value with the vessel FULL; the engine scales it with the wetted wall, UA(h) = UA$\cdot$A\_w(h)/A\_w(H), A\_w = $\pi$Dh + $\pi$D$^2$/4 (166.7 W/K at the starting 0.5 m, announced). The reaction heat r$\cdot$V follows the volume too. 4. The ledger claims mass and energy and closes (\textasciitilde{}1e-10 or better): both species carry a formation datum, so the vessel takes the canonical energy route, which the run announces.}
\tutreadme{check\_dynamic\_cstr\_level re-runs this case and holds these numbers.}
\begin{tutgolden}
\goldrow{kpi}{FeedStep}{MV\_final}{0.0015}
\goldrow{kpi}{FeedStep}{PV\_final}{0}
\goldrow{kpi}{FeedStep}{SP\_final}{0.0015}
\goldrow{kpi}{LC1}{MV\_final}{0.00150002945263}
\goldrow{kpi}{LC1}{PV\_final}{0.5000019311}
\goldrow{kpi}{LC1}{SP\_final}{0.5}
\goldrow{kpi}{TC1}{MV\_final}{386.663508487}
\goldrow{kpi}{TC1}{PV\_final}{339.991752771}
\goldrow{kpi}{TC1}{SP\_final}{340}
\goldrow{kpi}{balance}{element\_C\_closure\_rel}{1.66272359835e-14}
\goldrow{kpi}{balance}{element\_H\_closure\_rel}{1.65023582724e-14}
\goldrow{kpi}{balance}{element\_O\_closure\_rel}{1.66272359835e-14}
\goldrow{kpi}{balance}{element\_worst\_closure\_rel}{1.66272359835e-14}
\goldrow{kpi}{balance}{energy\_balance\_available}{1}
\goldmore{27}
\end{tutgolden}

\clearpage
\tutentry{ctrl23\_fedbatch\_yeast\_mazzoleni2015}{tutorials/ctrl/ctrl23\_fedbatch\_yeast\_mazzoleni2015}
\begin{tutscheme}
\begin{tikzpicture}[>={Stealth[length=2mm]}, every node/.style={font=\scriptsize}]
  \node[draw, rounded corners=2pt, align=center, minimum width=1.9cm, minimum height=0.85cm, fill=black!4] (reactor) at (0.0,-0.0) {\textbf{reactor}\\[-1pt]{\tiny dynamicCSTR}};
  \draw[->] ($(reactor.west)+(-1.05,0.00)$) -- ($(reactor.west)+(0,0.00)$) node[at start, above, font=\tiny, anchor=south west] {feed};
  \draw[->] ($(reactor.east)+(0,0.00)$) -- ($(reactor.east)+(1.05,0.00)$) node[at end, above, font=\tiny, anchor=south east] {broth};
\end{tikzpicture}
\end{tutscheme}
\tutwhat{Fed-batch of S. cerevisiae CEN.PK113-7D on glucose, the seven balances and the rate laws of Mazzoleni et al. 2015 (Table 4 parameters, nothing fitted): 17 h of batch, then an exponential glucose feed of 0.16 1/h to 49 h in a variable-volume vessel}
\begin{tutnarrative}
an aerobic fed-batch of yeast, from a published model with every parameter printed.

SOURCE (read in full, CC BY 4.0): Mazzoleni S, Landi C, Carten? F, de Alteriis E, Giannino F, Paciello L, Parascandola P, "A novel process-based model of microbial growth: self-inhibition in Saccharomyces cerevisiae aerobic fed-batch cultures", Microb. Cell Fact. 14 (2015) 109, doi:10.1186/s12934-015-0295-4. READ: Eq 1-7 (the seven mass balances), Table 1 (the rate laws), Table 2 (the auxiliary equations), Table 3 (initial values and feed set-up) and Table 4 (all parameters, four strains), and the Methods paragraphs on the CEN.PK cultures. NOT READ: the figures as data (Figs 2-7 are plots only), the paper's supplementary files, the optimisation of the calibration.

WHICH RUN. CEN.PK113-7D, the prototroph of Fig 4: feed 500 g/l glucose, mu = 0.16 1/h, feeding from t\_F = 17 h to t\_END = 49 h, the M\_F and y\_R of its Table 3. Why this one: it is the run whose biomass reaches the highest density (about 100 g/l), so the inhibitor feedback that is the paper's point acts; it is a new experiment of the paper's own, with the three measured variables of its figure; and the concentrated feed keeps the volume change small. The paper ran the strain with its own set of Table 4 values, used here without change.

THE MODEL. Seven pools in grams: G glucose, P glycolysis products, E ethanol, C\_M carbon metabolites (the cell material, B = P + C\_M the active mass), I inhibitor, R reserves, D dead cells. dG/dt = Feeding - Uptake dP/dt = eta\_G Uptake - Resp\_P - Ferm - Accum - Death\_P dE/dt = eta\_FE Ferm - Resp\_E dC\_M/dt = eta\_RP Resp\_P + eta\_RE Resp\_E + eta\_FP Ferm - Secretion - Death\_CM dR/dt = eta\_A Accum - Death\_R dI/dt = Secretion dD/dt = Death\_P + Death\_CM + Death\_R The laws are in constant/reactions, one reaction per process (eight), each an `expression` law with every Table 4 value and its table reference beside it.

HOW THE ENGINE CARRIES IT. (1) A `dynamicCSTR` with `volumeMode variable` and a zero outflow IS a fed-batch vessel: it fills and its volume follows the inventory, so every c[i] = n\_i / V is the paper's [G] = G/V. (2) Each pool is a mass-unit pseudo-component (MW = 1, volume \textasciitilde{} 0), so kmol are kg and kmol/m3 are g/l; the water is real and carries the volume. (3) The paper's efficiencies eta < 1 lose mass; the lost part is the named component `lostMass` (about 14 to 20 \% of each flux), so the mass balance closes and the case shows what the paper leaves untold. (4) The exponential feed F0 exp(mu (t - t\_F)) of Table 2 is a `Signal` staircase on the vessel's inflow, 64 steps of half an hour whose value is the EXACT MEAN of the exponential over the step, so the volume fed at every step boundary is exact (the feed rate within a step is flat: this is the discretisation, 0.16 1/h x 0.5 h = 8 \% between steps). (5) The vessel is adiabatic and all pool records share one formation datum: no metabolic heat (the paper has none), so T stays at 303.15 K.

WHAT THE PAPER'S TABLE 2 GIVES FOR V(t) is not used as printed. It reads V = M\_F/(c\_F y\_R) exp(mu (t - t\_F)) , a few ml at t\_F for this run, the volume of FEED that has entered when F0 = M\_F mu/(c\_F y\_R) is integrated -- but the vessel held 1 l (Methods) at t\_F, which is V\_0. The case uses the mass-conserving form V = V\_0 + (M\_F/(c\_F y\_R)) (exp(mu (t - t\_F)) - 1), which is what the staircase integrates to, and reproduces the paper's reported \textasciitilde{}100 g/l at 40 h.

WHAT IS NOT MODELLED, from the paper itself: oxygen, pH, the base added for pH control, antifoam and the volume they add; the evaporation of ethanol; the sampling. The paper's data are NOT reproduced as a pinned anchor: Fig 4 gives mass, glucose and ethanol as curves, and the comparison in README.md is information, read off the figure by eye.

\end{tutnarrative}
\tutreadme{An aerobic fed-batch of $\ast$Saccharomyces cerevisiae$\ast$ CEN.PK113-7D on glucose: 17 h of batch on 20 g/l, then an exponential feed of 500 g/l glucose at 0.16 1/h until 49 h. The model is the one of Mazzoleni et al. 2015 ($\ast$Microb. Cell Fact.$\ast$ 14:109, doi 10.1186/s12934-015-0295-4, CC BY 4.0): seven mass balances, eight rate laws, every parameter printed. Nothing is fitted here: the numbers are the paper's Table 3 and Table 4 (CEN.PK113-7D column). Run it: runCase tutorials/ctrl/ctrl23\_fedbatch\_yeast\_mazzoleni2015 (it runs under choupoCtrl, about a minute and a half).}
\tutreadme{READ, in the full text: Eq 1-7 (the seven balances), Table 1 (rate laws), Table 2 (auxiliary equations), Table 3 (initial values and feed set-up), Table 4 (all parameters, four strains), the Methods on the CEN.PK cultures (1 l of initial medium, 2 \% w/v glucose, 50 \% w/v glucose feed, pH 5, DOT 30 \%) and Fig 4 as an image. NOT read: the data behind the figures (only plots exist), the paper's other strains' runs, the supplementary material.}
\tutreadme{Why this run: CEN.PK113-7D (Fig 4) reaches the highest cell density of the paper (about 100 g/l), where the inhibitor feedback it is about acts, and it is the paper's own new experiment with mass, glucose and ethanol measured.}
\tutreadme{batchReactor could not carry this: it has no continuous feed and its volume is a number it keeps while a recipe transfer adds liquid, so every c = n/V would be wrong. The variable-volume dynamicCSTR under choupoCtrl is the engine's fed-batch vessel. No engine code was added.}
\tutreadme{Table 2 prints V(t) = M\_F/(c\_F y\_R) exp(?(t - t\_F)), which is the volume of FEED entered (14.6 ml at t\_F), not the volume of the vessel (1 l by the Methods). The case uses V = V$_0$ + (M\_F/(c\_F y\_R))(exp(?(t - t\_F)) - 1), which is what the staircase integrates to; it is declared in the header. Two other readings are declared in constant/reactions: ? (secretion) is printed in 1/h but is a dimensionless fraction, and the death switch d is written as a ramp of 0.001 g/l around $\tau$ (the expression grammar has no comparison). The temperature (303.15 K) is an assumption: the paper states 30 $^\circ$C only for its plate counts; it affects only the density of water, because the model is isothermal and the pseudo-components share one formation datum.}
\tutreadme{1. The diauxic batch. Glucose is gone at 14 h, ethanol peaks at about 9.5 g/l near 12.5 h and is consumed by 21 h (trajectory.csv). 2. The feed. After 17 h glucose stays near 0.01 g/l: the yeast takes what is fed and the biomass follows the feed, up to about 90 g/l at 40.5 h. 3. The overflow. After 40 h the glycolysis products pile up, ethanol is made and glucose accumulates; dead cells appear after 41 h. At 49 h the vessel holds 183.6 g of ethanol, 493.7 g of glucose and 60.9 g of dead cells in 3.436 l. 4. The mass balance closes (1.6e-14 relative; the energy ledger 3e-14). The paper does not track the mass its efficiencies lose; here it is 340.8 g of lostMass at 49 h, 27.5 \% of all the solute, so the case shows what the model leaves untold. The solutes in the vessel equal the 20 g charged plus the glucose fed (500 g/l $\times$ the fed volume): 1.2376 kg against 1.2382 kg (0.05 \%), the difference being the volume the pseudo-components are given (1e-9 m$^3$/mol, a total of 1.2 ml).}
\begin{tutgolden}
\goldrow{kpi}{ExpFeed}{MV\_final}{5.810215e-05}
\goldrow{kpi}{ExpFeed}{PV\_final}{0}
\goldrow{kpi}{ExpFeed}{SP\_final}{5.810215e-05}
\goldrow{kpi}{balance}{energy\_balance\_available}{1}
\goldrow{kpi}{balance}{energy\_closure\_rel}{3.25753563777e-14}
\goldrow{kpi}{balance}{energy\_functional\_H}{1}
\goldrow{kpi}{balance}{energy\_kJ\_final}{-54275.9280895}
\goldrow{kpi}{balance}{energy\_kJ\_heat\_cum}{0}
\goldrow{kpi}{balance}{energy\_kJ\_in\_cum}{-38468.6530934}
\goldrow{kpi}{balance}{energy\_kJ\_initial}{-15807.2749961}
\goldrow{kpi}{balance}{energy\_kJ\_out\_cum}{0}
\goldrow{kpi}{balance}{energy\_residual\_kJ}{-1.76805770025e-09}
\goldrow{kpi}{balance}{mass\_closure\_rel}{1.55410854116e-14}
\goldrow{kpi}{balance}{mass\_kg\_final}{4.6577532578}
\goldmore{31}
\end{tutgolden}

